% !TeX program = xelatex
% Build: bash slides/build.sh  (= mkdir -p build + two xelatex passes; four SimplePlus .sty files bundled)
% Manual: xelatex -interaction=nonstopmode -halt-on-error -output-directory=build mse_slides.tex  (twice)
% Content: ads-proj1 迷你搜索引擎项目报告（reports/，2026-10-20 终稿）的幻灯版。
% 全部数字来自 reports/data/ 证据链（results-tasks.md / results-verify.md），
% 与 REVIEW.md 数字审计一致；外推点均在页面上标注"外推"。
% 按课程要求匿名：不含姓名/学号/个人信息。
\documentclass[UTF8,fontset=none,aspectratio=169,10pt]{ctexbeamer}
\usepackage{amsmath,amssymb,booktabs,tabularx,array}
\usepackage{pgfplots}
\pgfplotsset{compat=1.18}
\usetheme{SimplePlus}
\usetikzlibrary{arrows.meta,positioning}
% 字体不再交给 ctex 的 fontset=windows：它按名字引用 SimSun/SimHei/KaiTi/微软雅黑，
% 系统没装就报 fontspec "cannot be found" 并让中文整页缺字。这里显式设置并逐级回落：
% 本机 Windows 字体 → Office 自带的 msyh.ttc → PingFang SC → Source Han Sans SC。
\newcommand{\SlidesCJK}[2]{%
  \setCJKmainfont{#1}[#2]%
  \setCJKsansfont{#1}[#2]%
  \setCJKmonofont{#1}[#2]%
  \setCJKfamilyfont{zhsong}{#1}[#2]%
  \setCJKfamilyfont{zhhei}{#1}[#2]%
  \setCJKfamilyfont{zhfs}{#1}[#2]%
  \setCJKfamilyfont{zhkai}{#1}[#2]%
  \setCJKfamilyfont{zhyahei}{#1}[#2]%
}
\IfFontExistsTF{Microsoft YaHei}
  {\SlidesCJK{Microsoft YaHei}{BoldFont=Microsoft YaHei Bold}}
  {\IfFileExists{/Applications/Microsoft Word.app/Contents/Resources/DFonts/msyh.ttc}
     {\SlidesCJK{msyh.ttc}{Path=/Applications/Microsoft Word.app/Contents/Resources/DFonts/,BoldFont=msyhbd.ttc}}
     {\IfFontExistsTF{PingFang SC}
        {\SlidesCJK{PingFang SC}{}}
        {\SlidesCJK{Source Han Sans SC}{}}}}
\colorlet{primary}{DarkBlue}
\colorlet{accent}{MediumBlue}
\definecolor{alertred}{RGB}{178,34,34}
\setbeamerfont{block title}{size=\normalsize,series=\mdseries}
\setbeamerfont{frametitle}{size=\large,series=\bfseries}
\setbeameroption{hide notes}
\setlength{\emergencystretch}{1.5em}
\newcolumntype{L}{>{\raggedright\arraybackslash}X}
\newcolumntype{R}{>{\raggedleft\arraybackslash}X}
\newcommand{\key}[1]{\textcolor{accent}{\textbf{#1}}}
\newcommand{\bad}[1]{\textcolor{alertred}{\textbf{#1}}}
\newcommand{\src}[1]{\par\vspace{0.5em}{\fontsize{8}{9.6}\selectfont\color{gray}#1\par}}
\newcommand{\gap}{\par\vspace{0.7em}}
\newcommand{\vsp}{\vspace*{\stretch{1}}}
\newcommand{\code}[1]{\texttt{\detokenize{#1}}}
\newcommand{\bigstat}[1]{{\Large\bfseries\color{accent}#1}}
% 标题页：去掉空白的 author 框，内容整体垂直居中（主题默认偏上，下半页留白过多）
\setbeamertemplate{title page}{%
  \vfill
  \begingroup
  \centering
  \begin{beamercolorbox}[sep=10pt,center]{title}
    \usebeamerfont{title}\inserttitle\par
    \vskip0.3em
    {\usebeamerfont{subtitle}\usebeamercolor[fg]{subtitle}\insertsubtitle\par}
  \end{beamercolorbox}
  \vskip0.9em
  \begin{beamercolorbox}[sep=6pt,center]{institute}
    \usebeamerfont{institute}\insertinstitute
  \end{beamercolorbox}
  \vskip0.4em
  \begin{beamercolorbox}[sep=6pt,center]{date}
    \usebeamerfont{date}\insertdate
  \end{beamercolorbox}
  \endgroup
  \vfill
}
% 目录条目间距（主题默认 4 行挤在一起）
\setbeamertemplate{section in toc}{%
  \leavevmode\inserttocsectionnumber.\ \inserttocsection\par\vspace{0.8em}}
% 表格统一放大：11pt 正文用 \normalsize 表格，列距略放宽
\setlength{\tabcolsep}{4pt}
\renewcommand{\arraystretch}{1.15}
\title[迷你搜索引擎]{迷你搜索引擎}
\subtitle{倒排索引的构建、查询与阈值分析\quad·\quad 数据结构课程项目 8-1}
\author{}
\institute{C++17\quad 主语料：莎士比亚全集（195 篇）\quad 副语料：IETF RFC ×3}
\date{2026 年 10 月}

\begin{document}

% ============================ F1 标题 ============================
\begin{frame}
  \titlepage
  \note{报告 21 页正文 + 完整证据链；本幻灯只引用证据编号（F 系列），全部数字与 reports/data/ 对账过。}
\end{frame}

% ============================ F2 概览 ============================
\begin{frame}{概览}
  \vsp
  \begin{columns}[T]
    \begin{column}{0.5\textwidth}
      \tableofcontents
    \end{column}
    \begin{column}{0.46\textwidth}
      \begin{block}{三组数字}
        \normalsize
        \begin{tabularx}{\linewidth}{L r}
          语料规模（主语料） & \bigstat{195 篇} \\
          正文词元 & \bigstat{866,626} \\
          正式索引词项 & \bigstat{14,937} \\
          正式索引体积 & \bigstat{3.69 MB} \\
          阈值实验测量 & \bigstat{201 条} \\
          CTest 断言 & \bigstat{85,155} \\
        \end{tabularx}
      \end{block}
    \end{column}
  \end{columns}
  \vsp
  \normalsize 四项课程任务全部完成；Bonus 完成。
  % \src{报告摘要、§3、§4；REVIEW.md 数字审计}
\end{frame}

\section{设计流程}

% ============================ F3 任务与语料 ============================
\begin{frame}{任务与语料}
  \begin{block}{课程任务：4 项 + Bonus}
    \small
    \begin{enumerate}\setlength{\itemsep}{0.15em}
      \item 词频统计与\key{停用词}界线
      \item 带\key{Porter 词干化}的倒排索引（不含停用词）
      \item 单词 / 短语（及布尔）查询，返回文档 ID
      \item 查询阈值实验：\code{max-df}、\code{min-df}、\code{min-match}、\code{top-x}
      \item[] \hspace{-1.2em}\key{Bonus}：\key{50 万文件} + \key{4 亿不同词}
    \end{enumerate}
  \end{block}
  \vsp
  \begin{center}
  \small
  \begin{tabularx}{0.95\textwidth}{L r r}
    \toprule
    语料 & 主语料：莎士比亚全集 & 副语料：IETF RFC \\
    \midrule
    文档数 & 195（37 剧 / 154 十四行诗 / 4 长诗） & 3 \\
    词数（空白切分） & 930,911 & 81,848 \\
    正文词元（默认配置） & 866,626 & 75,250 \\
    语料字节 & 4,919,371 B & 569,472 B \\
    \bottomrule
  \end{tabularx}
  \end{center}
  \vspace{0.4em}
  % \src{tasks F1/F3；verify F10；报告摘要、§4.1}
\end{frame}

% ============================ F4 流水线 ============================
\begin{frame}{构建流水线}
  \vsp
  \begin{center}
    \begin{tikzpicture}[
      nd/.style={draw=MediumBlue, fill=MutedBlue, rounded corners=3pt, font=\small,
                 inner sep=5pt, minimum height=8ex, minimum width=2.7cm, align=center},
      gen/.style={draw=MediumGreen, fill=MutedGreen, rounded corners=3pt, font=\small,
                  inner sep=5pt, minimum height=8ex, minimum width=2.7cm, align=center},
      ar/.style={-{Stealth[length=2.6mm]}, draw=MediumBlue, thick},
      node distance=4pt]
      \node[nd] (a) {语料镜像\\与抽取\\{\scriptsize 1,211→195 篇}};
      \node[nd, right=of a] (b) {分词\\~\\{\scriptsize 866,626 词元}};
      \node[nd, right=of b] (c) {停用词\\过滤\\{\scriptsize 移除 35.74\%}};
      \node[nd, right=of c] (d) {Porter\\词干化\\{\scriptsize 25,705 → 14,947}};
      \node[gen, right=of d] (e) {倒排索引\\v1\\{\scriptsize 3,685,655 B}};
      \draw[ar] (a) -- (b);
      \draw[ar] (b) -- (c);
      \draw[ar] (c) -- (d);
      \draw[ar] (d) -- (e);
    \end{tikzpicture}
  \end{center}
  \vsp
  \begin{columns}[T]
    \begin{column}{0.48\textwidth}
      \begin{block}{构建侧}
        \small
        \begin{itemize}\setlength{\itemsep}{0.1em}
          \item 词典可插拔：伸展树 / 哈希 / \code{std::map}
          \item 索引可写 v1（原始 u32）或 v2（差分 + varint）
          \item 大语料走 SPIMI 式块式构建
        \end{itemize}
      \end{block}
    \end{column}
    \begin{column}{0.48\textwidth}
      \begin{block}{查询侧}
        \small
        \begin{itemize}\setlength{\itemsep}{0.1em}
          \item 单词 / 短语 / 布尔，输出 JSON
          \item 四类阈值：\code{max-df}、\code{min-df}、\code{min-match}、\code{top-x}
          \item 冷查询走磁盘 B+ 树附属词典
        \end{itemize}
      \end{block}
    \end{column}
  \end{columns}
  \vsp
  % \src{报告 §1–§2；tab:coverage}
\end{frame}

% ============================ F5 词频与 Zipf ============================
\begin{frame}{词频统计：Zipf 头部集中}
  \begin{columns}[T]
    \begin{column}{0.44\textwidth}
      \begin{block}{词表结构（默认配置）}
        \small
        \begin{tabularx}{\linewidth}{L r}
          不同词形 & 25,705 \\
          仅出现一次（hapax） & 9,753 \\
          top-10 词元占比 & 20.41\% \\
          top-100 词元占比 & 54.30\% \\
          top-1,000 词元占比 & 80.26\% \\
        \end{tabularx}
      \end{block}
    \end{column}
    \begin{column}{0.54\textwidth}
      \centering
      \begin{tikzpicture}
        \begin{axis}[
          width=6.6cm, height=4.2cm,
          xmode=log, ymode=log,
          xlabel={词频排名 $r$}, ylabel={词频 tf},
          axis lines=left, axis line style={gray!60},
          tick label style={font=\scriptsize},
          label style={font=\small},
          grid=major, grid style={gray!20},
          every axis plot/.append style={line width=0pt},
          title style={font=\small},
        ]
          \addplot[only marks, mark=*, mark size=0.9pt, color=accent]
            table[x=rank,y=tf] {../reports/figures/data/zipf.dat};
        \end{axis}
      \end{tikzpicture}
      \par{\scriptsize\color{gray}主语料 tf–排名双对数图：近似 $1/r$ 衰减（图 1）}
    \end{column}
  \end{columns}
  \vsp
  \begin{block}{观察}
    \footnotesize
    极少数词覆盖大量词元，这正是停用词值得单独处理的统计原因；
    同时 hapax 占词表 \key{37.9\%}，是词干化能大幅合并词表的空间所在。
  \end{block}
  % \src{tasks F6–F8；wordcount-default.csv；报告 §2.4}
\end{frame}

% ============================ F6 停用词 ============================
\begin{frame}{停用词：从人工候选带到统计规则}
  \begin{block}{第一阶段：tf/df 候选带扫描 + 人工定界}
    \normalsize
    按候选带 $4 \to 9 \to 14 \to 16 \to 21 \to 31 \to 48$ 词逐步放开，观察 tf/df 曲线拐点，
    人工定界得到 \key{39 词}。
  \end{block}
  \vsp
  \begin{block}{第二阶段：``宽、重、平''三特征规则（纯统计，无人为挑选）}
    \normalsize
    \begin{tabularx}{\linewidth}{L L}
      \key{宽}：$\mathrm{df} \ge 0.40 \times$ 文档数 &
      \key{重}：tf 占比 $\ge 0.001$ \\
      \key{平}：长度归一化速率熵 $\ge 0.93$ &
      \key{结果}：同样 39 词，与人工表重叠 35 \\
    \end{tabularx}
  \end{block}
  \vsp
  \begin{block}{效果与边界}
    \normalsize
    39 个词只占词表的 \key{0.15\%}，却覆盖 \key{35.74\%} 的正文词元（移除 309,760 个）。
    速率熵度量``分布是否平坦''，是三特征里区分功能词与内容词的决定性一条。
  \end{block}
  % \src{tasks F9/F11–F14；报告 §2.5、§3.5}
\end{frame}

% ============================ F7 RFC 交叉验证 ============================
\begin{frame}{停用词没有万能表：RFC 交叉验证}
  \begin{center}
  \small
  \begin{tabularx}{0.9\textwidth}{L r r}
    \toprule
    同一套三特征规则 $(0.40,\,0.001,\,0.93)$ & 莎士比亚 & IETF RFC \\
    \midrule
    推导出的停用词数 & 39 & \key{58} \\
    与莎士比亚 39 词表的重叠 & — & 仅 26 词 \\
    索引体积（原始 → 正式配置） & 5,598,793 → 3,685,655 B & 510,181 → 308,164 B \\
    \bottomrule
  \end{tabularx}
  \end{center}
  \vspace{0.4em}
  \small
  \begin{itemize}\setlength{\itemsep}{0.35em}
    \item RFC 2119 关键词 \code{must} / \code{may} / \code{should}：df 比值均为 1（足够``宽''``重''），
      但速率熵 \key{0.809 / 0.899 / 0.885} $< 0.93$（不够``平''）→ \key{正确排除}
    \item \code{server}、\code{page}、\code{rfc} 速率熵接近 1 → 收入推导表
      （\code{[Page N]} 样板属 producer 侧噪声，非引擎问题）
    \item 端到端：RFC 上原始配置 75,250 词元 → 正式配置 45,370 词元
  \end{itemize}
  \vsp
  \begin{block}{结论}
    \small 同一套阈值在不同语料上给出不同的停用词表——停用词是\key{语料相关}的，
    规则可以复用，但停用词表需要重新选取。
  \end{block}
  % \src{verify F10；报告 §4.1、表 tab:rfc}
\end{frame}

% ============================ F8 词干化与索引配置 ============================
\begin{frame}{Porter 词干化与四种索引配置}
  \begin{block}{自实现 Porter 词干化}
    \small
    25,705 词形 → \key{14,947 词干}（\key{−41.85\%}），hapax 9,753 → 4,843；
    词干化把``同源词形''折叠为一个词项，直接缩小词表与倒排记录。
  \end{block}
  \vspace{0.6em}
  \begin{center}
  \small
  \begin{tabularx}{0.95\textwidth}{L r r r}
    \toprule
    配置（flags） & 词项数 & 倒排记录 & 索引字节 \\
    \midrule
    raw（无过滤、无词干） & 25,705 & 152,650 & 5,598,793 \\
    stop（停用词） & 25,666 & 147,909 & 4,320,622 \\
    stem（词干化） & 14,947 & 121,903 & 4,958,598 \\
    \key{正式（stopwords + stemmed）} & \key{14,937} & \key{117,704} & \key{3,685,655} \\
    \bottomrule
  \end{tabularx}
  \end{center}
  \vsp
  \begin{block}{结论}
    \small 停用词与词干化分别压缩``词元流长度''与``词表大小''，两者叠加使正式配置最小。
  \end{block}
  % \src{tasks F18；报告表 tab:index-configs}
\end{frame}

% ============================ F9 索引格式 v1/v2 ============================
\begin{frame}{索引格式：v1 布局与 v2 压缩}
  \begin{columns}[T]
    \begin{column}{0.47\textwidth}
      \begin{block}{v1：原始 u32 布局}
        \small
        \begin{itemize}\setlength{\itemsep}{0.1em}
          \item 文件头 + 文档表 + 有序词表 + 倒排记录区
          \item 倒排记录：docID、tf、位置，各 4 字节
        \end{itemize}
      \end{block}
    \end{column}
    \begin{column}{0.49\textwidth}
      \begin{block}{v2：只改倒排记录区}
        \small
        \begin{itemize}\setlength{\itemsep}{0.1em}
          \item docID 与位置存\key{差分}
          \item tf 与差分值用\key{字节对齐变长码}（varint）
        \end{itemize}
      \end{block}
    \end{column}
  \end{columns}
  \vsp
  \begin{center}
  \small
  \begin{tabularx}{0.95\textwidth}{L r r r r}
    \toprule
    语料 & 第一版字节 & 第二版字节 & 第二版/第一版 & 压缩构建峰值内存 \\
    \midrule
    主语料（莎士比亚） & 3,685,655 & \key{1,725,694} & \key{46.82\%} & 0.19 s（内建） \\
    500k 合成语料 & 427,603,247 & 284,560,009 & 66.55\% & 1.46 GiB（v1 为 1.01） \\
    \bottomrule
  \end{tabularx}
  \end{center}
  \vsp
  \begin{block}{成本权衡}
    \footnotesize
    \begin{itemize}\setlength{\itemsep}{0.1em}
      \item 主语料倒排记录区占 86.0\%、段内节省 61.8\% → 整体节省 53.18\%
      \item 500k 词项表占比高（倒排区仅 56.1\%），段内节省 59.6\% 与主语料接近
      \item 压缩不是全维度占优：构建峰值内存 1.01 → \bad{1.46 GiB}；查询输出与 v1 完全一致
    \end{itemize}
  \end{block}
  % \src{tasks F21；verify F13/F14；报告 §3.4}
\end{frame}

% ============================ F10 词典对比 ============================
\begin{frame}{三种词典：构建过程不同，产物相同}
  \vsp
  \begin{center}
  \normalsize
  \begin{tabularx}{0.95\textwidth}{L r r r}
    \toprule
    词典 & 比较 / 旋转 & 额外事件 & 构建耗时 \\
    \midrule
    伸展树（自实现） & 6,392,724 / 5,850,795 & — & 0.19 s \\
    哈希表 & 688,494 & 11 次 rehash & \key{0.09 s} \\
    \code{std::map}（基准参考） & 9,435,694 & — & 0.14 s \\
    \bottomrule
  \end{tabularx}
  \end{center}
  \vsp
  \begin{block}{结论}
    \normalsize
    哈希最快；伸展树最慢，但中序遍历\key{直接产出有序词表}。
    三种词典写出的索引\key{相同}，说明词典只改变构建过程，
    不改变索引结果。
  \end{block}
  % \src{tasks F20；verify F5；报告 §4.3}
\end{frame}

% ============================ F11 B+ 树 sidecar + SPIMI ============================
\begin{frame}{磁盘 B+ 树附属词典与 SPIMI 块式构建}
  \begin{columns}[T]
    \begin{column}{0.48\textwidth}
      \begin{block}{B+ 树 sidecar（查询侧）}
        \normalsize
        词项 → 倒排记录（偏移，块长）的磁盘映射：
        \begin{itemize}\setlength{\itemsep}{0.25em}
          \item 文件仅 \key{237,568 B}（索引的 6.45\%）
          \item 冷查询只读 \key{8,308 B} = 全量加载的 \key{0.225\%}
          \item 与全量加载逐词一致（双路径校验）
        \end{itemize}
      \end{block}
    \end{column}
    \begin{column}{0.48\textwidth}
      \begin{block}{SPIMI 块式构建（构建侧）}
        \normalsize
        内存预算 4,096 活词：
        \begin{itemize}\setlength{\itemsep}{0.25em}
          \item 17 个临时块、峰值活词 \key{5,129}
          \item 临时块共 4,166,556 B
          \item 归并产物与内存构建\key{一致}
        \end{itemize}
      \end{block}
    \end{column}
  \end{columns}
  \vsp
  \begin{block}{为什么是这两个部件}
    \normalsize
    两者合起来去掉``构建/查询必须把全词表放进内存''的假设：
    查询侧把词典移到磁盘，构建侧把内存压到 $O(\text{预算})$ 与词表规模无关——
    这是 Bonus 两条规模轴的工程基础。
  \end{block}
  % \src{tasks F22；verify F23；报告 §2.7、§2.10}
\end{frame}

\section{测试与结果}

% ============================ F12 查询 ============================
\begin{frame}{查询设计与样例}
  \small
  支持单词、短语（词元流上整段匹配）与布尔（AND 语义，按命中词数排序）；输出 JSON，\code{--pretty} 可读。
  \vspace{0.4em}
  \begin{center}
  \begin{tabularx}{0.98\textwidth}{>{\raggedright\arraybackslash}p{3.4cm} L r}
    \toprule
    查询 & 行为 & 结果 \\
    \midrule
    \code{osric} & 稀有词，单文档角色 & 1 篇 \\
    \code{love} & 高频词（top-3 tf 195 / 165 / 142） & 138 篇 \\
    \code{"king richard"} & 短语整段匹配 & 7 篇（top-3 tf 193 / 191 / 140） \\
    \code{"my lord"} & \code{my} 为停用词 → 只查 \code{lord} & 42 篇 \\
    \code{king queen} & 布尔：按命中词数 + tf 排序 & 33 篇双命中 / 43 篇单命中 \\
    \bottomrule
  \end{tabularx}
  \end{center}
  \vsp
  \begin{block}{短语语义的边界}
    \small
    位置定义在\key{过滤后}的词元流上：\code{"my lord"} 退化为 \code{lord}，
    全由停用词组成的短语返回空——这是设计选择，不是实现缺陷。
  \end{block}
  \vsp
  \small 排序仅用命中词数与 tf 总和，\bad{无 idf 与长度归一化}，长剧更易靠前。
  % \src{tasks F27–F33；报告 §3.8}
\end{frame}

% ============================ F13 阈值实验 ============================
\begin{frame}{阈值实验：阈值改变的是查询语义}
  \normalsize
  15 条查询 $\times$ 四类阈值 = \key{201 条测量}（中位 19.2 ms，总计 4.05 s）；
  用 Jaccard 系数对比阈值前后的结果集。
  \vsp
  \begin{block}{反直觉但可解释的三个现象}
    \normalsize
    \begin{enumerate}\setlength{\itemsep}{0.3em}
      \item \bad{收紧 $\ne$ 变少}：\code{sweet sword moon} 加 \code{--min-df 50} 后 \code{sword} 等词被裁，
        布尔查询退化为单词 \code{moon}，结果 \key{30 → 90 篇}（Jaccard 0.333）
      \item \key{短语不受词侧阈值影响}：\code{"king richard"} 全档 7 → 7（Jaccard 1.00）——
        短语必须整段匹配，裁剪查询词只会让短语失效或不变
      \item \key{停用词查询恒为空}：\code{the} / \code{of} 全档 0 → 0——过滤索引的设计行为，非 bug
    \end{enumerate}
  \end{block}
  \vsp
  \begin{columns}[T]
    \begin{column}{0.48\textwidth}
      \small\code{--top-x} 才是真正的\key{结果截断}：
      love 的 top-5 / 10 / 20 对应 Jaccard 0.036 / 0.073 / 0.145。
    \end{column}
    \begin{column}{0.48\textwidth}
      \small\code{--min-match} 改变\key{布尔条件}（AND 要求命中词数）：
      1 / 2 / 3 对应 93 / 39 / 30 篇。
    \end{column}
  \end{columns}
  \vsp
  \footnotesize\key{一句话总结}：\code{max-df}/\code{min-df} 是查询词筛选器（偏精度），
  \code{min-match} 改布尔语义，\code{top-x} 才截断结果——四者作用在不同阶段，
  所以``阈值收紧''不能简单等同于``结果变少''。
  % \src{tasks F36–F39；threshold\_matrix.csv；报告 §3.9}
\end{frame}

\section{Bonus 的结果与分析}

% ============================ F14 B+ 树规模 ============================
\begin{frame}{Bonus 1：磁盘 B+ 树词典的规模实测}
  \normalsize
  4,908,747 词项的内存索引峰值 2.00 GiB → 推测 4 亿词 $\approx$ \bad{160 GiB}，
  单机放不下：词典必须上磁盘。
  \vsp
  \begin{center}
  \small
  \begin{tabular}{r r r r r r}
    \toprule
    键数 & 树高 & 总页数 & 文件字节 & 字节/键 & 页读/查询 \\
    \midrule
    100,000 & 2 & 369 & 1,515,520 & 15.1552 & 2 \\
    1,000,000 & 3 & 3,688 & 15,110,144 & 15.1101 & 3 \\
    10,000,000 & 3 & 36,865 & 151,003,136 & 15.1003 & 3 \\
    30,000,000 & 3 & 110,593 & 452,993,024 & 15.0998 & 3 \\
    50,000,000 & 4 & 184,322 & 754,987,008 & 15.0997 & 4 \\
    \bottomrule
  \end{tabular}
  \end{center}
  \vsp
  \begin{block}{外推与边界}
    \small
    4 KiB 页、批量装载、每档 1,000 次随机查询全部校验；体积近似\key{严格线性}（15.10 B/键）。
    外推 4 亿键 $\approx$ \key{6.04 GB}、树高 \key{4}、单查询读 4 页 = \key{16 KiB}
    （树高与页读已在 5,000 万键\key{实测}到 4，外推的只有文件字节与页数）。
  \end{block}
  % \src{verify F11；bptree-bench.csv；报告表 tab:bptree-scale}
\end{frame}

% ============================ F15 500k 实验 ============================
\begin{frame}{Bonus 2：500,000 文档合成语料（实测）}
  \normalsize
  合成语料：500,000 个文件、20,000,000 词元、4,908,747 个不同词、19,999,929 条倒排记录。
  \vsp
  \begin{center}
  \normalsize
  \begin{tabularx}{0.98\textwidth}{L r r r}
    \toprule
    构建方式 & 耗时 & 峰值内存 & 产物 \\
    \midrule
    块式 v1（4,855 块） & 扫描 5.47 + 归并 125.08 s & \key{1.01 GiB} & 427,603,247 B \\
    块式 v2（差分压缩） & 扫描 5.42 + 归并 123.31 s & 1.46 GiB & 284,560,009 B \\
    内存构建 v1 & 总 76.75 s（内建 44.93 s） & 2.00 GiB & 427,603,247 B \\
    \bottomrule
  \end{tabularx}
  \end{center}
  \vsp
  \begin{block}{结论}
    \normalsize
    块式与内存构建的 v1 索引\key{逐字节相同}（SHA-256 一致）：内存 $O(\text{预算})$ 与词表规模无关，
    在 50 万文件规模\key{实测成立}。代价是约 \bad{2.2 倍墙钟}——
    归并现为 $O(T\cdot R)$ 线性扫游标，换成堆可降到 $O(T\log R)$。
  \end{block}
  % \src{verify F12–F15；报告表 tab:500k}
\end{frame}

% ============================ F16 正确性 ============================
\begin{frame}{正确性：可复现的验证链}
  \normalsize
  \begin{itemize}\setlength{\itemsep}{0.5em}
    \item CTest \key{18/18} 通过：8 个 Catch2 可执行文件、46 个用例、\key{85,155 条断言}（B+ 树占 99.3\%）
    \item 使用了 Python 进行\key{对比检验}：
      \begin{itemize}\setlength{\itemsep}{0.2em}
        \item 四种配置词频 CSV 逐字节一致（1,677,253 / 1,674,689 / 958,633 / 957,871 B）
        \item 正式 / 原始 / 压缩索引的全词表 df、tf、docID、位置全项一致
        \item 短语对拍：305 条采样、303 条检查、\key{0 不一致}（21.20 s）
        \item 磁盘词典双路径：B+ 树查询与全量加载逐词一致
      \end{itemize}
    \item 讲义四文档黄金样例逐项吻合；块式 / 内存 / 三词典产物 SHA-256 相同
  \end{itemize}
  \vsp
  \begin{block}{项目的最大价值}
    \normalsize
    不是某个墙钟数字，而是这条验证链：同一文本层能被独立实现重建，三种实现写出同一字节流，
    ``看起来能跑''的错误在提交前就被暴露。
  \end{block}
  % \src{verify F1–F9；tasks F15；报告 §3.1–§3.4、§4.4}
\end{frame}

\section{项目评价}

% ============================ F17 局限与结论 ============================
\begin{frame}{局限与结论}
  \begin{columns}[T]
    \begin{column}{0.48\textwidth}
      \begin{block}{主要局限}
        \small
        \begin{itemize}\setlength{\itemsep}{0.25em}
          \item 排序无 idf / BM25，长剧占优
          \item 短语匹配在过滤后词元流上（\code{"my lord"} 退化）
          \item 切词仅 ASCII；构建与查询单线程
          \item 无增量索引；4 亿词未端到端实测
        \end{itemize}
      \end{block}
    \end{column}
    \begin{column}{0.48\textwidth}
      \begin{block}{下一步（按收益排序）}
        \small
        \begin{enumerate}\setlength{\itemsep}{0.25em}
          \item 粒度实验（scene 级 / 合并短诗）
          \item BM25 / TF-IDF 与相关性评价
          \item SPIMI 归并输出直接写 B+ 树，跑一次真正 500k/400M 端到端
        \end{enumerate}
      \end{block}
    \end{column}
  \end{columns}
  \vsp
  \begin{block}{结论}
    \small
    四项课程任务全部完成并经独立对拍；\key{50 万文件轴实测通过}，
    \key{4 亿词轴}给出实测锚点上的外推（6.04 GB / 树高 4 / 4 页读 = 16 KiB），
    两条轴在全文都明确区分``实测''与``外推''。
  \end{block}
  \vsp
  \centering\small\color{gray}完整停用词表见附录 A / 附录 B。
  % \src{报告 §3.5、§4.2；REVIEW.md 数字审计}
\end{frame}

% ============================ F18 附录 A：主语料停用词 ============================
\begin{frame}{附录 A：主语料的停用词（人工 39 / 统计 39）}
  \begin{columns}[T]
    \begin{column}{0.49\textwidth}
      \begin{block}{人工定界（39 词，\code{data/stopwords.txt}）}
        \footnotesize
        the, and, i, to, of, a, my, that, in, is, not, me, it, for, with, be,
        his, this, but, have, as, so, what, by, no, all, do, if, are, on, from,
        or, which, more, then, their, when, than, yet
      \end{block}
    \end{column}
    \begin{column}{0.49\textwidth}
      \begin{block}{三特征统计推导（39 词，df$\ge$0.40、tf\_share$\ge$0.001、熵$\ge$0.93）}
        \footnotesize
        the, and, i, to, of, a, my, that, in, is, not, me, it, for, with, be,
        his, this, but, have, as, \key{thou}, so, what, by, no, all, do, if,
        are, \key{thee}, on, from, \key{o}, which, then, when, \key{hath}, yet
      \end{block}
    \end{column}
  \end{columns}
  \vspace{0.5em}
  \begin{block}{两份表的差异：重叠 35 / 39}
    \small
    推导\key{多出} thou、thee、o、hath（古英语功能词——``分布平坦''能自动识别它们）；
    人工表\key{多出} more、their、than、or（比较与选择连词，速率熵略低于 0.93，被判为``不够平''）。
  \end{block}
  % \src{tasks F9/F11–F14；data/stopwords.txt；derived-df0.40-er0.93.txt；logs/tasks/stopwords-list.log}
\end{frame}

% ============================ F19 附录 B：RFC 停用词 ============================
\begin{frame}{附录 B：副语料 RFC 的停用词（58 词，纯统计推导）}
  \centering
  \footnotesize
  \begin{tabular*}{0.97\textwidth}{@{\extracolsep{\fill}}lllll}
    the & a & to & of & and \\
    in & is & that & for & be \\
    \key{rfc} & or & not & \key{server} & as \\
    this & are & an & by & it \\
    with & if & \key{data} & from & \key{page} \\
    \key{information} & on & used & when & any \\
    standards & use & can & other & one \\
    e & will & some & have & all \\
    \key{protocol} & no & only & case & its \\
    more & \key{list} & \key{error} & \key{time} & received \\
    each & does & same & they & g \\
    \key{transfer} & however & multiple & & \\
  \end{tabular*}
  \vspace{0.15em}
  \begin{block}{读法}
    \footnotesize
    同一套规则 $(0.40,\,0.001,\,0.93)$ 得到 58 词：\key{加粗}的 server / page / rfc / data /
    information / protocol / error / time / list / transfer 是典型\key{领域内容词}，
    而 RFC 2119 的 must / may / should（熵 0.809 / 0.899 / 0.885 $<$ 0.93）\bad{不在表中}。
  \end{block}
  \vsp
  \centering\small\color{gray}谢谢！问题与讨论欢迎提出。
  % \src{verify F10；work/verify/rfc-stopwords.txt（e / g 来自 ``e.g.'' 切分）；报告 §4.1}
\end{frame}

\end{document}
